\documentclass[12pt]{article}

\usepackage[T1]{fontenc}
\usepackage{amsthm}
\usepackage{amsmath}
\usepackage{amsfonts}
\usepackage[scr]{rsfso}
\usepackage{bm}
\usepackage{thmtools}
\usepackage{mathtools}
\usepackage{enumitem}
\usepackage{hyperref}
\usepackage{cleveref}
\usepackage{cancel}
\usepackage{booktabs}
\usepackage{graphicx}
\usepackage{subcaption}
\usepackage{booktabs}
\usepackage{siunitx}

\usepackage{natbib}
\bibliographystyle{plainnat}


\usepackage[parfill]{parskip}
\setlength{\parindent}{0pt}
\setlength{\parskip}{2ex plus 0.2ex minus 0.2ex}

\usepackage[
    paper  = letterpaper,
    left   = 1in,
    right  = 1in,
    top    = 1in,
    bottom = 1in,
]{geometry}



\newcommand{\N}{\mathbb{N}}
\newcommand{\Z}{\mathbb{Z}}
\newcommand{\Q}{\mathbb{Q}}
\newcommand{\R}{\mathbb{R}}
\newcommand{\C}{\mathbb{C}}

\newcommand{\GG}{\mathcal{G}}
\newcommand{\LL}{\mathcal{L}}
\newcommand{\NN}{\mathcal{N}}

\newcommand{\iid}{\overset{iid}{\sim}}
\renewcommand{\Pr}{\mathbb{P}}
\newcommand{\E}{\mathbb{E}}

\newcommand{\ELBO}{\text{ELBO}}

\title{
    Image Segmentation with Hierarchical Potts-Markov Priors
}
\author{
    Satvik Saha\footnote{\texttt{ss7228@columbia.edu}}\\
    \textit{\normalsize Department of Statistics, Columbia University}
}
\date{}

\begin{document}

\maketitle


\begin{abstract}
    Image segmentation generally describes the process of separating an image into
    different groups or regions of pixels, with pixels in a group sharing
    similar characteristics.
    A crude method of image segmentation involves clustering pixels by their color,
    with no regard to their spatial arrangement.
    In this work, we introduce spatial structure to a mixture model for
    segmenting images, by imposing suitable priors on group labels.
    In particular, we focus on hierarchical Potts-Markov-like priors which
    encourage smoothing based on pixel neighborhoods and group-specific
    temperatures.
    We use MCMC-based inference procedures to apply this method on the Berkeley
    Segmentation Dataset, and find that it generally outperforms the
    corresponding spatially-na\"ive mixtures.
\end{abstract}



\section{Introduction}


Image segmentation, the process of dividing an image into different parts
(segments), is a problem in computer vision that has applications in many
areas, ranging from biomedical imaging (detecting diseased tissue) to
autonomous vehicles (detecting objects on the road).
Ideally, each segment represents a recognizable object in a scene, i.e.\ is a
region of similar pixels.
This lets us frame segmentation as a clustering problem, for which mixture
models are well suited.

Interpreting an image as a bag of pixel colors $\bm{x} = \{x_\ell\}_{\ell \in
\LL}$ suggests a crude segmentation method: pick $K$ colors, then declare a
pixel to be part of the $k$-th object if its color is most similar to the
$k$-th color.
While such a procedure works reasonably well, it often leads to noisy segments
with fuzzy boundaries; this kind of algorithm discards a lot of potentially
useful information, in particular the spatial configuration of the pixels in
the image.
For instance, supposing that our objects of interest occupy continuous regions
in space, we may believe that neighboring pixels will more often than not
belong to the same segment.
In this project, we build a simple, spatially-aware probabilistic model that
exploits this heuristic and produces cleaner segmentations than a na\"ive
mixture.

\begin{figure}[ht]
\begin{center}
    \begin{subfigure}[t]{0.32\textwidth}
    \begin{center}
        \includegraphics[width=\textwidth]{figures/42049_orig}
    \end{center}
    \caption{Original image}
    \label{fig:42049_orig}
    \end{subfigure}
    %
    \begin{subfigure}[t]{0.32\textwidth}
    \begin{center}
        \includegraphics[width=\textwidth]{figures/42049_seg_human}
    \end{center}
    \caption{Human segmentation}
    \label{fig:42049_human}
    \end{subfigure}
    %
    \begin{subfigure}[t]{0.32\textwidth}
    \begin{center}
        \includegraphics[width=\textwidth]{figures/42049_K4_seg_kmeans}
    \end{center}
    \caption{$K$-Means}
    \label{fig:42049_kmeans}
    \end{subfigure}

    \vspace{0.5em}

    \begin{subfigure}[t]{0.32\textwidth}
    \begin{center}
        \includegraphics[width=\textwidth]{figures/42049_K4_seg_gmm}
    \end{center}
    \caption{GMM}
    \label{fig:42049_gmm}
    \end{subfigure}
    %
    % \begin{subfigure}[t]{0.32\textwidth}
    % \begin{center}
    %     \includegraphics[width=\textwidth]{figures/42049_K4_seg_fixed}
    % \end{center}
    % \caption{GMM-PM-($\bm{\beta} = \bm{1}$)}
    % \label{fig:42049_fixed}
    % \end{subfigure}
    %
    \begin{subfigure}[t]{0.32\textwidth}
    \begin{center}
        \includegraphics[width=\textwidth]{figures/42049_K4_seg}
    \end{center}
    \caption{GMM-PM \eqref{eq:potts_mixture}}
    \label{fig:42049_potts}
    \end{subfigure}
    %
    \begin{subfigure}[t]{0.32\textwidth}
    \begin{center}
        \includegraphics[width=\textwidth]{figures/42049_K4_beta}
    \end{center}
    \caption{Posterior summary of $\bm{\beta}$}
    \label{fig:42049_beta}
    \end{subfigure}
\end{center}
\caption{%
    Segmentations of an image (\ref{fig:42049_orig}, \texttt{42049} in the
    Berkeley Segmentation Dataset) by humans (\ref{fig:42049_human}), $K$-Means
    (\ref{fig:42049_kmeans}), a simple Gaussian Mixture (GMM,
    \ref{fig:42049_gmm}), and a Gaussian Mixture with a hierarchical
    Potts-Markov prior (GMM-PM, \ref{fig:42049_potts}).
    The last three models use $K = 4$ segments.
    The posterior distribution of $\bm\beta$ in GMM-PM (\ref{fig:42049_beta})
    gives us information about the relative smoothness of different segments.
}
\label{fig:img_42049}
\end{figure}



\subsection{Image Segmentation with Mixture Models}

A general mixture model for $\bm{x} = \{x_\ell\}_{\ell \in \LL}$ looks like
\begin{equation}
\label{eq:general_mixture}
\begin{aligned}
    \bm{\beta}
        &\sim p(\bm{\beta}), \\
    \bm{\theta}
        &\sim p(\bm{\theta}), \\
    \bm{z} \mid \bm\beta
        &\sim p(\bm{z} \mid \bm\beta), \\
    x_\ell \mid \bm\theta, z_\ell = k
        &\sim p(x_\ell \mid \bm\theta, z_\ell = k) =: f(x_\ell; \theta_k),
            &\quad\ell \in \LL.
\end{aligned}
\end{equation}
Each data item $x_\ell$ is associated with a latent label $z_\ell \in \{1,
\dots, K\}$ which determines which component it belongs to.
The distribution $f(\,\cdot\,; \theta_k)$ of items within the $k$-th component
is controlled by the parameter $\bm\theta$.

A simple mixture model may choose \emph{a priori} independent labels $z_\ell
\mid \bm\beta$, i.e.\ allow \begin{equation}
\label{eq:naive_prior}
    p(\bm{z} \mid \bm\beta) = \prod_{\ell \in \LL} p(z_\ell \mid \bm\beta)
\end{equation}
to factorize.
This makes the pairs $\{(z_\ell, x_\ell)\}_{\ell \in \LL}$
exchangeable: the location $\ell \in \LL$ of a data point $x_\ell$ plays no
role in the model.
In the context of image segmentation, with $x_\ell \in \R^d$ denoting the color
of a pixel, its location $\ell$ within the image is ignored.

Instead, we will choose our prior $p(\bm{z}\mid\bm\beta)$ in a way that
respects the spatial structure of $\LL$.
Specifically, we would like to use the following heuristic ---
\emph{neighboring pixels often belong to the same segment} --- motivated by the
following observation.
A segment within an image is rarely uniform in color; small fluctuations and
scattered islands of pixels with very different colors may appear within the
larger region.
A simple mixture model using \eqref{eq:naive_prior} would see these islands as
part of a completely different segment (or as individual segments of their
own).
For example, looking closely at the segmentation produced by such a model in
Figure~\ref{fig:42049_gmm} reveals a visually noisy blue sky.
The kind of spatially aware model we explore here chooses to smooth out these
irregularities and group them as part of the encompassing segment.
This ought to result in a cleaner, less noisy segmented image; for instance,
note the more uniform sky in Figure~\ref{fig:42049_potts}.
Additionally, such a model can also give us the opportunity to learn about the
smoothness of each segment.

Note that this is very much a modeling choice/goal, which is not always
desirable: in applications like medical imaging, such small irregularities are
often the target and should \emph{not} be glossed over!
The aim here is to identify high level structure in an image where all objects
are of similar importance, rather than picking out fine detail.



\subsection{The Potts-Markov Prior}
\label{sec:potts_prior}

\begin{figure}
\begin{center}
    \begin{subfigure}[t]{0.24\textwidth}
    \begin{center}
        \includegraphics[width=\textwidth]{figures/potts_beta_0}
    \end{center}
    \caption{$\bm\beta = \bm0$}
    \label{fig:potts_beta_0}
    \end{subfigure}
    %
    \begin{subfigure}[t]{0.24\textwidth}
    \begin{center}
        \includegraphics[width=\textwidth]{figures/potts_beta_1}
    \end{center}
    \caption{$\bm\beta = \bm1$}
    \label{fig:potts_beta_1}
    \end{subfigure}
    %
    \begin{subfigure}[t]{0.24\textwidth}
    \begin{center}
        \includegraphics[width=\textwidth]{figures/potts_beta_10}
    \end{center}
    \caption{$\bm\beta = \bm{10}$}
    \label{fig:potts_beta_10}
    \end{subfigure}
    %
    \begin{subfigure}[t]{0.24\textwidth}
    \begin{center}
        \includegraphics[width=\textwidth]{figures/potts_beta_custom}
    \end{center}
    \caption{$\bm\beta = (1, 1.5, 1.5)^\top$}
    \label{fig:potts_beta_custom}
    \end{subfigure}
\end{center}
\caption{%
    Draws from $\text{Potts}(\bm\beta)$, where each pixel is connected to its
    (at most) 8 immediate neighbors.
    Larger values of $\beta_k$ promote more cohesive components with smoother
    boundaries: components in \ref{fig:potts_beta_1} have noticeably more
    `islands' of foreign labels compared to \ref{fig:potts_beta_10}.
    Note that imbalanced $\bm\beta$, as in \ref{fig:potts_beta_custom}, leads
    to some components effectively vanishing.
}
\label{fig:potts_samples}
\end{figure}



Let $\GG(\LL, \mathcal{E})$ be a graph on the set of locations $\LL$, and
notate $\ell \sim \ell'$ when $\ell, \ell' \in \LL$ are neighbors in $\GG$.
We denote $\bm{z} \sim \text{Potts}(\bm\beta)$ for $\bm\beta \succeq \bm0$,
when the distribution of $\bm{z}$ has density
\begin{equation}
\label{eq:potts}
    p(\bm{z} \mid \bm\beta) = \frac{1}{Z(\bm{\beta})} \exp\left(\sum_{\ell \sim \ell'} \sum_{k = 1}^K \beta_k \bm{1}(z_{\ell'} = z_{\ell} = k)\right).
\end{equation}
This produces a smoothing effect on the class labels $\bm{z}$ with respect to
the spatial structure $\GG$.
The `reward' for agreement between neighboring nodes assigned class $k$ is
controlled by a class-specific temperature $\beta_k$: this allows different
components to have smoother boundaries than others.
See \Cref{fig:potts_samples} for illustrative samples from
$\text{Potts}(\bm\beta)$.
% By further imposing a fairly non-informative prior on $\bm\beta$, its posterior
% distribution can reveal the relative smoothness of each segment.

Observe that \eqref{eq:potts} is an exponential family in $\bm\beta$, with
sufficient statistic $\bm{T}(\bm{z})$ described by
\begin{equation}
    T_k(\bm{z})
        := \sum_{\ell\sim\ell'} \bm{1}(z_\ell = z_{\ell'} = k), \qquad\qquad 1 \leq k \leq K.
\end{equation}
Each $T_k$ counts the number of edges in $\GG$ connecting items
assigned the same class $k$.
With this, we have $p(\bm{z}\mid\bm{\beta}) = \exp(\bm\beta^\top
\bm{T}(\bm{z})) / Z(\bm\beta)$.
Further keeping track of the statistics
\begin{equation}
    U_{k\ell}(\bm{z}) := \sum_{\ell \sim \ell'} \bm{1}(z_{\ell'} = k), \quad\qquad 1 \leq k \leq K, \quad \ell \in \LL,
\end{equation}
which count the number of neighbors of $z_\ell$ assigned the class $k$, the
full conditionals
\begin{equation}
    p(z_{\ell} = k\mid \bm{z}_{-\ell}, \bm\beta) \propto \exp\left(\beta_k U_{k\ell}(\bm{z})\right)
\end{equation}
suggest a simple Gibbs sampling algorithm for $\text{Potts}(\bm\beta)$.


% \emph{Remark:} Note that \[
%     T_k(\bm{z})
%         = \frac{1}{2}\sum_{\ell \in \LL} U_{k\ell}(\bm{z}) \,\bm{1}(z_\ell = k).
% \]

The partition function $Z(\bm\beta) = \sum_{\bm{z} \in \{1, \dots, K\}^{|\LL|}}
\exp(\bm\beta^\top \bm{T}(\bm{z}))$ is an intractable normalizing constant.
This makes \emph{both} the joint density of the mixture model
\eqref{eq:general_mixture}, as well as the full-conditional
\begin{equation}
\label{eq:beta_full_conditional}
    p(\bm\beta \mid \bm{z}, \bm{\theta}, \bm{x})
        \propto \frac{p(\bm\beta)}{Z(\bm\beta)} \exp(\bm\beta^\top \bm{T}(\bm{z}))
\end{equation}
intractable, which presents problems for standard posterior inference
algorithms such as Gibbs sampling and mean field variational inference.
For the former choice, sampling from \eqref{eq:beta_full_conditional} using
ordinary Metropolis-Hastings requires computing $Z(\bm\beta^*)/Z(\bm\beta)$ as
part of the acceptance ratio.
This issue of a \emph{doubly intractable likelihood} can be sidestepped using
an \emph{auxiliary variable} technique.



\subsection{The Exchange Algorithm}
\label{sec:exchange_algorithm}

The following auxiliary variable Metropolis-Hastings scheme \citep{murray2006}
is used to produce samples from \eqref{eq:beta_full_conditional} without
directly dealing with the partition function $Z(\,\cdot\,)$.
\begin{enumerate}[itemsep=-0.7em]
    \item Propose $\bm{\beta}^* \sim q(\bm{\beta}^*; \bm{\beta})$.
    \item Draw auxiliary $\bm{z}^* \sim p(\bm{z}^* \mid \bm{\beta}^*)$.
    \item Accept $\bm{\beta}^*$ with probability
    \begin{equation}
    \label{eq:exchange_ratio}
        \frac{
            q(\bm{\beta}; \bm{\beta}^*)\,
            p(\bm{\beta}^*)
        }{
            q(\bm{\beta}^*; \bm{\beta})\,
            p(\bm{\beta})
        }\cdot \frac{
            p(\bm{z} \mid \bm{\beta}^*)
        }{
            p(\bm{z} \mid \bm{\beta})
        }\cdot \frac{
            p(\bm{z}^*\mid \bm{\beta})
        }{
            p(\bm{z}^* \mid \bm{\beta}^*)
        },
    \end{equation}
    and repeat from step 1.
\end{enumerate}
Note that the problematic terms $Z(\bm\beta)$, $Z(\bm\beta^*)$ neatly cancel
out in the ratio \eqref{eq:exchange_ratio}.
This is a simpler version of the algorithm in \citet{moller2006}.
% An easy way to see why this algorithm works is to observe that this is simply
% Metropolis-Hastings for drawing samples of $(\bm{\beta}', \bm{z}')$ from
% $p(\bm\beta'\mid \bm{z})\, p(\bm{z}' \mid \bm{\beta}')$, with proposal
% $q(\bm\beta'^*; \bm\beta')\, p(\bm{z}'^* \mid \bm\beta'^*)$; the marginal for
% $\bm\beta'$ is our target.

Running this algorithm is fairly expensive in practice for image data,
especially since step 2 involves sampling from $\text{Potts}(\bm{\beta}^*)$,
with the benefit of targeting the correct stationary distribution.
Mean field variational methods are much faster, at the cost of using inexact
approximations for $Z(\,\cdot\,)$.
See Section~\ref{sec:inference}, step (c) for details on implementation.




\section{Gaussian Mixture Model with Potts-Markov Prior}
\label{sec:model}

Let $x_\ell \in \R^d$ be the color of the pixel at location $\ell \in \LL$.
We use a $K \in \N$ component mixture model of the form
\eqref{eq:general_mixture}, with a Potts-Markov prior \eqref{eq:potts} on the
latent class labels $\bm{z} = \{z_\ell\}_{\ell \in \LL}$, multivariate
Gaussians as component distributions with means and covariances $\{(\mu_k,
\Sigma_k)\}_{k = 1}^K \subseteq \R^d \times \R^{d\times d}$, and conjugate
Normal-Inverse-Wishart priors on the latter component parameters.
We further impose weak Gamma priors on the smoothness parameters $\bm\beta$,
which yields the following hierarchical model.
% Let $\LL$ be the set of locations of pixels, and for $\ell \in
% \LL$, let $\NN_\ell \subseteq \LL$ be the set of
% neighbors of $\ell$ in a graph $\GG(\LL, \mathcal{E})$.
% Denote $\ell \sim \ell'$ when $\ell, \ell'$ are neighbors in $\GG$,
% i.e.\ $\{\ell, \ell'\} \in \mathcal{E}$.
% Let $x_\ell \in \R^d$ be the intensity of the pixel at $\ell$.
% We model $\bm{x} := \{\bm{x}\}_{\ell \in \LL}$ as a mixture of $K \in
% \N$ Gaussians, with class means and covariances $\{(\mu_k, \Sigma_k)\}_{k =
% 1}^K \subseteq \R^d \times \R^{d\times d}$.
% Each pixel $x_\ell$ is associated with a class label $z_\ell \in \{1, \dots,
% K\}$.
\begin{equation}
\label{eq:potts_mixture}
\begin{aligned}
    \beta_k
        &\iid \text{Gamma}(a_\beta, b_\beta),
        & 1 \leq k \leq K, \\
    \bm{z} \mid \bm{\beta}
        &\sim \text{Potts}(\bm\beta), \\
    \Sigma_k
        &\iid \text{Inverse-Wishart}(\Sigma_0, \nu_0),
        & 1 \leq k \leq K, \\
    \mu_k \mid \Sigma_k
        &\sim \text{Normal}(m_0, \kappa_0^{-1}\Sigma_k),
        & 1 \leq k \leq K, \\
    x_\ell \mid \bm{\mu}, \bm{\Sigma}, z_{\ell} = k
        &\sim \text{Normal}(\mu_k, \Sigma_k),
        & \ell \in \LL.
\end{aligned}
\end{equation}
We will abbreviate \eqref{eq:potts_mixture} as a GMM-PM model.
Its joint density factors as
\begin{equation}
\label{eq:joint_density}
    \prod_{k = 1}^K p(\beta_k)
        \cdot p(\bm{z}\mid\bm{\beta})
        \cdot \prod_{k = 1}^K p(\Sigma_k)\, p(\mu_k \mid \Sigma_k)
        \cdot \prod_{\ell \in \LL} p(x_\ell \mid \mu_k, \Sigma_k, z_\ell = k).
\end{equation}

Here, $\LL = \{(i, j): 1 \leq i \leq m, 1 \leq j \leq n\}$ is a
collection of tuples representing positions on an $m\times n$ image grid, and
we choose the spatial structure $\GG$ such that each pixel is connected to its
(at most) 8 immediate neighbors on the grid.
Formally, the neighborhood of $\ell$ in $\GG$ looks like $\NN_{\ell} = (\ell +
\Delta) \cap \LL$ for the set of translations $\Delta = \{-1, 0,
1\}^2\setminus\{(0, 0)\}$.
Furthermore, we have $d = 3$, with the components of $x_\ell$ denoting red,
green, blue intensities in $[0, 1]$.

Modifying \eqref{eq:potts_mixture} by fixing $\bm\beta = \bm0$ reduces it to a
simple Gaussian mixture model with \emph{a priori} independent labels $\bm{z}$.
In contrast, the Potts-Markov prior on $\bm{z} \mid \bm{\beta}$ induces the
full conditionals
\begin{equation}
    p(z_\ell = k \mid \bm{z}_{-\ell}, \beta_k, \mu_k, \Sigma_k, x_\ell)
        \propto \exp\left(\beta_k U_{k\ell}(\bm{z})\right) \cdot p(x_\ell \mid \mu_k, \Sigma_k, z_\ell = k),
\end{equation}
where the weights $\exp(\beta_k U_{k\ell})$ encourage each $z_\ell$ to conform
with its neighbors.
% This has a smoothing effect on the class labels $\bm{z}$ with respect to the
% spatial structure of the image.
% The `reward' for agreement is controlled by class-specific temperatures
% $\beta_k$, which allows different classes to have smoother boundaries than
% others.
% The partition function $Z(\bm{\beta})$ for $p(\bm{z} \mid \bm{\beta})$ is in
% general intractable, which makes sampling from the full-conditional of
% $\bm{\beta}$ difficult in a Gibbs sampling inference procedure; this can be
% circumvented using an auxiliary variable Metropolis-Hastings technique.

% Define the statistics \[
%     T_k(\bm{z})
%         := \sum_{\ell\sim\ell'} \bm{1}(z_\ell = z_{\ell'} = k), \qquad\qquad 1 \leq k \leq K
% \] which count the number of edges connecting pixels assigned the same class
% $k$.
% Then, we may alternatively describe \[
%     p(\bm{z}\mid\bm{\beta})
%         = \frac{1}{Z(\bm{\beta})} \exp\left(\sum_{k = 1}^K \beta_k T_k(\bm{z}) \right)
%         = \frac{1}{Z(\bm{\beta})} \exp\left(\bm{\beta}^\top \bm{T}(\bm{z})\right).
% \]
% Note that $\bm{T}(\bm{z})$ is a sufficient statistic for the exponential family
% $p(\bm{z} \mid \bm{\beta})$.

% Further keeping track of the statistics \[
%     U_{k\ell}(\bm{z}) = \sum_{\ell \sim \ell'} \bm{1}(z_{\ell'} = k)
% \] which count the number of neighbors of $z_\ell$ assigned the class $k$, note
% that the full conditional \[
%     p(z_{\ell} = k\mid \bm{z}_{-\ell}, \bm{\beta}, \bm{\mu}, \bm{\Sigma}, \bm{x})
%         \propto \exp\left(\beta_k U_{k\ell}\right) \cdot p(x_\ell \mid \mu_k, \Sigma_k, z_\ell = k).
% \]


% Note that \[
%     p(\bm{\beta} \mid \bm{z})
%         = \frac{p(\bm\beta)}{Z(\bm\beta)}\, \exp\left(\bm{\beta}^\top \bm{T}(\bm{z})\right)
% \] only depends on $\bm{z}$ through $\bm{T}(\bm{z})$.



\section{Data and Implementation}
\label{sec:data_implementation}

\begin{figure}[ht!]
\begin{center}
    \begin{subfigure}[t]{0.32\textwidth}
    \begin{center}
        \includegraphics[width=\textwidth]{figures/108073_orig}
    \end{center}
    \caption{Original image}
    \label{fig:108073_orig}
    \end{subfigure}
    %
    \begin{subfigure}[t]{0.32\textwidth}
    \begin{center}
        \includegraphics[width=\textwidth]{figures/108073_seg_human}
    \end{center}
    \caption{Human segmentation}
    \label{fig:108073_human}
    \end{subfigure}
    %
    \begin{subfigure}[t]{0.32\textwidth}
    \begin{center}
        \includegraphics[width=\textwidth]{figures/108073_K3_seg_kmeans}
    \end{center}
    \caption{$K$-Means}
    \label{fig:108073_kmeans}
    \end{subfigure}

    \vspace{0.5em}

    \begin{subfigure}[t]{0.32\textwidth}
    \begin{center}
        \includegraphics[width=\textwidth]{figures/108073_K3_seg_gmm}
    \end{center}
    \caption{GMM}
    \label{fig:108073_gmm}
    \end{subfigure}
    %
    % \begin{subfigure}[t]{0.32\textwidth}
    % \begin{center}
    %     \includegraphics[width=\textwidth]{figures/108073_K3_seg_fixed}
    % \end{center}
    % \caption{GMM-PM-($\bm{\beta} = \bm{1}$)}
    % \label{fig:108073_fixed}
    % \end{subfigure}
    %
    \begin{subfigure}[t]{0.32\textwidth}
    \begin{center}
        \includegraphics[width=\textwidth]{figures/108073_K3_seg}
    \end{center}
    \caption{GMM-PM \eqref{eq:potts_mixture}}
    \label{fig:108073_potts}
    \end{subfigure}
    %
    \begin{subfigure}[t]{0.32\textwidth}
    \begin{center}
        \includegraphics[width=\textwidth]{figures/108073_K3_beta}
    \end{center}
    \caption{Posterior summary of $\bm{\beta}$}
    \label{fig:108073_beta}
    \end{subfigure}
\end{center}
\caption{%
    Segmentations of an image (\ref{fig:108073_orig}, \texttt{108073} in the
    Berkeley Segmentation Dataset), with
    \labelcref{fig:108073_kmeans,fig:108073_gmm,fig:108073_potts} using $K = 3$
    segments.
    While $K$-Means and the GMM retain a lot of detail (the tiger's stripes,
    ripples in the water), GMM-PM attempts to smooth them out in favor of
    high-level structure.
}
\label{fig:img_108073}
\end{figure}



\begin{figure}[ht!]
\begin{center}
    \begin{subfigure}[t]{0.32\textwidth}
    \begin{center}
        \includegraphics[width=\textwidth]{figures/86000_orig}
    \end{center}
    \caption{Original image}
    \label{fig:86000_orig}
    \end{subfigure}
    %
    \begin{subfigure}[t]{0.32\textwidth}
    \begin{center}
        \includegraphics[width=\textwidth]{figures/86000_seg_human}
    \end{center}
    \caption{Human segmentation}
    \label{fig:86000_human}
    \end{subfigure}
    %
    \begin{subfigure}[t]{0.32\textwidth}
    \begin{center}
        \includegraphics[width=\textwidth]{figures/86000_K5_seg_kmeans}
    \end{center}
    \caption{$K$-Means}
    \label{fig:86000_kmeans}
    \end{subfigure}

    \vspace{0.5em}

    \begin{subfigure}[t]{0.32\textwidth}
    \begin{center}
        \includegraphics[width=\textwidth]{figures/86000_K5_seg_gmm}
    \end{center}
    \caption{GMM}
    \label{fig:86000_gmm}
    \end{subfigure}
    %
    % \begin{subfigure}[t]{0.32\textwidth}
    % \begin{center}
    %     \includegraphics[width=\textwidth]{figures/86000_K5_seg_fixed}
    % \end{center}
    % \caption{GMM-PM-($\bm{\beta} = \bm{1}$)}
    % \label{fig:86000_fixed}
    % \end{subfigure}
    %
    \begin{subfigure}[t]{0.32\textwidth}
    \begin{center}
        \includegraphics[width=\textwidth]{figures/86000_K5_seg}
    \end{center}
    \caption{GMM-PM \eqref{eq:potts_mixture}}
    \label{fig:86000_potts}
    \end{subfigure}
    %
    \begin{subfigure}[t]{0.32\textwidth}
    \begin{center}
        \includegraphics[width=\textwidth]{figures/86000_K5_beta}
    \end{center}
    \caption{Posterior summary of $\bm{\beta}$}
    \label{fig:86000_beta}
    \end{subfigure}
\end{center}
\caption{%
    Segmentations of an image (\ref{fig:86000_orig}, \texttt{86000} in the
    Berkeley Segmentation Dataset), with
    \labelcref{fig:86000_kmeans,fig:86000_gmm,fig:86000_potts} using $K = 5$
    segments.
    While $K$-Means and GMM retain more detail, the GMM-PM smooths them out and
    focuses on large-scale features.
    The posterior distribution of $\bm\beta$ (\ref{fig:86000_beta}) seems to
    indicate that the building and sky segments are smoother than the flower
    segments.
}
\label{fig:img_86000}
\end{figure}


We select color images from the Berkeley Segmentation Dataset
\citep{martin2001} and compare the segmentations produced by our GMM-PM model
\eqref{eq:potts_mixture} with those from $K$-Means algorithms and simple
Gaussian mixture models.
All images contain $|\LL| =$ 154,401 pixels (size $321\times 481$ or $481
\times 321$); the resulting spatial graph $\GG$ has $|\mathcal{E}| =$ 615,200
edges.

We use a Gibbs sampler, incorporating the exchange algorithm
(\ref{sec:exchange_algorithm}) to perform posterior inference on
GMM-PM \eqref{eq:potts_mixture}; given the sheer volume of data, most
hyperparameters are chosen such that the priors are non-informative.
See Section~\ref{sec:inference} for complete details on posterior inference,
Section~\ref{sec:hyperparameter} for hyperparameter selection.
The Gibbs sampler for GMM-PM has been implemented in \texttt{C++} with an
\texttt{R} wrapper (using \texttt{Rcpp}).
% The Gaussian mixture model is run using the same sampler, fixing $\bm\beta =
% \bm0$.
The $K$-means algorithm is run using a native R implementation.



\section{Evaluation and Results}

The dataset contains segmentations performed by humans, which we use as a
reference to compute Adjusted Rand Index (ARI) scores \citep{rand1971} for the
models of interest; see Table~\ref{tab:ari}.
These scores are a measure of similarity between two segmentations (not
necessarily having the same number of clusters): $1$ indicates a perfect match,
$0$ indicates a completely random segmentation.
\Cref{fig:img_42049,fig:img_108073,fig:img_86000} illustrate segmentations of
particular images.


\sisetup{detect-all=true}
\begin{table}[h!]
    \centering
    \caption{%
        Adjusted Rand Index (ARI) scores of segmentations against human
        references; higher is better.
        The last four columns report scores for the following models:
        $K$-Means, simple Gaussian Mixture (GMM), model
        \eqref{eq:potts_mixture} with fixed $\bm\beta = \bm{1}$
        (GMM-PM-($\bm{\beta} = \bm{1}$)), and model \eqref{eq:potts_mixture}
        (GMM-PM).
    }
    \label{tab:ari}
    \begin{tabular}{>{\tt}lc|SSSS}
    \toprule
    \normalfont{Image ID} & $K$
        & \text{$K$-Means}
        & \text{GMM}
        & \text{GMM-PM-($\bm\beta = \bm1$)}
        & \text{GMM-PM}
        \\\midrule
        108073 & 3
            &   0.2395
            &   0.2241
            &   \bfseries 0.2767
            &   0.2678
            \\
        198054 & 7
            & 0.3603
            & 0.4081
            & \bfseries 0.6521
            & 0.4444
            \\
        253027 & 3
            & 0.1707
            & 0.3470
            & \bfseries 0.4481
            & 0.4408
            \\
        253036 & 5
            & 0.4154
            & \bfseries 0.6150
            & 0.4444
            & 0.4459
            \\
        295087 & 8
            & 0.2591
            & 0.2693
            & 0.2967
            & \bfseries 0.2982
            \\
        300091 & 5
            & 0.1563
            & 0.1404
            & 0.1892
            & \bfseries 0.2648
            \\
        37073 & 5
            & 0.7151
            & 0.4705
            & 0.5428
            & \bfseries 0.5757
            \\
        38092 & 7
            & 0.4471
            & 0.5342
            & \bfseries 0.5453
            & 0.5381
            \\
        42049 & 4
            & 0.3009
            & 0.3187
            & \bfseries 0.3412
            & 0.3403
            \\
        86000 & 5
            & 0.4587
            & 0.4913
            & \bfseries 0.6341
            & 0.6270
            \\
    \bottomrule
    \end{tabular}
\end{table}

Table~\ref{tab:ari} indicates that although GMM-PM with fixed $\bm{\beta}$ can
outperform GMM-PM in full generality, the discrepancy is rarely large, and the
former loses out on information from the posterior distribution of the
smoothness parameter $\bm{\beta}$.
In general, human segmentations identify many more segments than our mixture
models, as they can tell apart different objects with the same color, which may
explain why the performances aren't particularly high.
To illustrate this point, note that in \Cref{fig:img_42049}, our models cannot
separate the bird and the branch which have very similar colors.


\section{Discussion}

We have presented a spatially aware Gaussian mixture model (GMM-MP), and
demonstrated that it frequently outperforms na\"ive Gaussian mixtures in the
context of image segmentation, especially when the focus is not on fine detail
but on broad strokes.
While the use of the Potts-Markov prior incorporates spatial information for
the purpose of smoothing (by rewarding cohesion), one can imagine priors which
encourage other kinds of behavior.
A generalization of \eqref{eq:potts} would be priors of the form
\begin{equation}
    p(\bm{z} \mid \bm{\beta})
        \propto \exp\left(\sum_{\ell \in \LL} \alpha(z_\ell; \ell) + \sum_{\ell, \ell' \in \LL} \beta(z_\ell, z_{\ell'}; \ell, \ell')\right)
\end{equation}
for potentials $\alpha, \beta$ that could model a variety of interactions
between labels.
Priors like these could be plugged directly into this project's setup with
little change to the core sampling algorithm.

Segmentation results such as those in
\Cref{fig:img_42049,fig:img_108073,fig:img_86000} show that a Gaussian mixture
is not always good at modeling colors that appear at the bundaries of the RGB
color space $[0, 1]^3$, including bright hilights and very dark areas.
This suggests that mixtures of Gaussians (which are supported on all of $\R^3$)
may not be the best choice for modeling color distributions.
We may swap out $\{f(\,\cdot\,; \theta)\}_{\theta \in \Theta}$ in
\eqref{eq:general_mixture} for any convenient family of likelihoods with
minimal change to the spatial prior.
Another solution may be to transform our color features $x_\ell \in [0, 1]^3$
to range over $\R^3$, say using logit links, before plugging them into our
model.


% \nocite{*}
\bibliography{references}




% \clearpage

\appendix




\section{Posterior Inference}
\label{sec:inference}

We produce posterior samples of $\bm{z}, \bm{\beta}, \bm{\mu}, \bm{\Sigma} \mid
\bm{x}$ using a Gibbs sampling scheme.
This involves updating a state $(\bm{z}, \bm{\beta}, \bm{\mu}, \bm{\Sigma})$ by
iteratively cycling through the following full-conditionals.
\begin{enumerate}[label=(\alph*)]
    \item Sample new component means and covariances $\bm{\mu}, \bm{\Sigma}
    \mid - $ as follows: for each $1 \leq k \leq K$, set \begin{align*}
        n_k &:= \sum_{\ell \in \LL} \bm{1}(z_\ell = k), &
        \bar{x}_k &:= \frac{1}{n_k}\sum_{\ell \in \LL} x_\ell \bm{1}(z_\ell = k), \\
        % \kappa_k &:= \kappa_0 + n_k, &
        % \nu_k &:= \nu_0 + n_k, \\
        \mu_k'&:= \frac{\kappa_0 m_0 + n_k \bar{x}_k}{\kappa_0 + n_k}, &
        S_k &:= \sum_{\ell \in \LL} (x_\ell - \bar{x}_k)(x_\ell - \bar{x}_k)^\top \bm{1}(z_\ell = k), \\
        \Sigma_k' &:= \Sigma_0 + S_k + \frac{\kappa_0 n_k}{\kappa_0 + n_k} (\bar{x}_k - m_0)(\bar{x}_k - m_0)^\top. \hspace{-8em}
    \end{align*}
    Then, replace the current $\bm{\mu}, \bm{\Sigma}$, by sampling from \begin{align*}
        \Sigma_k \mid -
            &\sim \text{Inverse-Wishart}\left(\Sigma_k', \nu_0 + n_k\right), \\\
        \mu_k \mid \Sigma_k, -
            &\sim \text{Normal}\left(\mu_k', (\kappa_0 + n_k)^{-1}\Sigma_k\right).
    \end{align*}
    These are standard Gaussian-Inverse-Wishart conjugate prior updates (see
    \cite{hoff2009}).

    \item Sample new labels $\bm{z} \mid -$ as follows: for each $\ell \in
    \LL$ (we choose a randomized order), sample $z_\ell$ from a
    categorical distribution with unnormalized probabilities \[
        p(z_\ell = k \mid -)
            \propto \exp\left(\beta_k U_{k\ell}\right)\cdot p(x_\ell \mid \mu_k, \Sigma_k, z_\ell = k)
    \] for $1 \leq k \leq K$.
    The second term is simply a (multivariate) Gaussian likelihood \[
        p(x_\ell \mid \mu_k, \Sigma_k, z_\ell = k)
            = (2\pi)^{-d/2}\det(\Sigma_k)^{-1 /2}\exp\left(-\frac{1}{2}(x_\ell - \mu_k)^\top \Sigma_k^{-1}(x_\ell - \mu_k)\right).
    \] These can be pre-computed and stored before looping over $\ell \in \LL$.
    The statistic $\bm{U}$ can be locally updated each time a new $z_\ell$ is
    sampled.

    \item Sample new Potts-Markov smoothness terms $\bm{\beta} \mid - $ using
    an auxiliary variable Metropolis-Hastings algorithm as follows: for $1 \leq
    t \leq T_{\beta}$, iterate the following.
    \begin{enumerate}[label=(\roman*)]
        \item Propose $\bm\beta^* \sim q(\,\cdot\,; \bm{\beta})$.
        Here, we use the proposal \[
            \log\bm{\beta}^* = \log\bm{\beta} + \bm\epsilon, \qquad
                \bm\epsilon \sim \text{Normal}\left(\bm0, \sigma_\beta^2 \bm{I}_K\right),
        \] to automatically keep $\bm\beta \succeq 0$.
        This makes the proposal ratio \[
            \frac{q(\bm{\beta}; \bm{\beta}^*)}{q(\bm{\beta}^*; \bm{\beta})}
                 = \prod_{k = 1}^K \exp(\epsilon_k).
        \]

        \item Draw an auxiliary set of labels $\bm{z}^* \sim p(\bm{z}^*\mid
        \bm{\beta}^*)$ from the Potts-Markov prior.
        This is accomplished by performing a few Gibbs sampling sweeps (with a
        hot-start initialization at the current $\bm{z}$), repeatedly looping
        through $\ell \in \LL$ and resampling $z_\ell^*$ with
        unnormalized probabilities \[
            p(z_\ell^* \mid - ) \propto \exp\left(\beta_k^* U_{k\ell}(\bm{z}^*)\right)
                =: \exp\left(\beta_k^* \sum_{\ell' \in \NN_\ell} \bm{1}(z_{\ell'}^* = k)\right).
        \]

        \item Update $\bm{\beta}$ to $\bm{\beta}^*$ with probability \[
            \frac{
                q(\bm{\beta}; \bm{\beta}^*)\,
                p(\bm{\beta}^*)
            }{
                q(\bm{\beta}^*; \bm{\beta})\,
                p(\bm{\beta})
            }\cdot \frac{
                p(\bm{z} \mid \bm{\beta}^*)
            }{
                p(\bm{z} \mid \bm{\beta})
            }\cdot \frac{
                p(\bm{z}^*\mid \bm{\beta})
            }{
                p(\bm{z}^* \mid \bm{\beta}^*)
            },
        \] keep the old value of $\bm{\beta}$ otherwise.
        Note that $p(\bm{\beta}), p(\bm{\beta}^*)$ are $\text{Gamma}(a_\beta,
        b_\beta)$ likelihoods, and the ratio \[
            \frac{
                p(\bm{z} \mid \bm{\beta}^*)
            }{
                p(\bm{z} \mid \bm{\beta})
            }\cdot \frac{
                p(\bm{z}^*\mid \bm{\beta})
            }{
                p(\bm{z}^* \mid \bm{\beta}^*)
            } = \exp\left(
                (\bm{\beta}^* - \bm{\beta})^\top (\bm{T}(\bm{z}) - \bm{T}(\bm{z}^*))
            \right)
        \] is free of the intractable partition function values $Z(\bm\beta),
        Z(\bm\beta^*)$!
    \end{enumerate}
\end{enumerate}

This describes a Gibbs within Metropolis-Hastings within Gibbs MCMC sampling
scheme.
In practice, we choose $5$ inner Gibbs loops (step (c)(ii)), $T_\beta = 10$
Metropolis-Hastings loops (step (c)), and $750$ outer Gibbs loops with the
first burn-in $250$ samples discarded.
% We also use a hot-start initialization of $(\bm{z}, \bm{\mu}, \bm{\Sigma})$
% using parameters obtained from a run of $K$-Means, while $\bm{\beta}$ is
% initialized to a vector of ones.
Convergence is assessed by examining trace plots of parameters.

One practical consideration is tuning the proposal step size $\sigma_\beta$.
This is adjusted dynamically during the burn-in period so that the acceptance
rate for $\bm\beta^*$ (step (c)) settles around 30\%.

We use the same sampling scheme with $\bm{\beta} = \bm0$ and skipping step (c)
to perform posterior inference on a Gaussian mixture model.


\section{Hyperparameter Selection}
\label{sec:hyperparameter}

Since the joint density \eqref{eq:joint_density} is intractable and there is no
convenient way of holding out image data once a spatial structure $\GG$ is
imposed, we estimate the required number of components $K$ by maximizing an
expected log predictive density score from a simpler Gaussian mixture (fixing
$\bm\beta = \bm0$ in \ref{eq:potts_mixture}), with the idea that this reduction
already models the color distribution of an image well and is thus a good
proxy.
The remaining hyperparameters in \eqref{eq:potts_mixture} are chosen to be
weakly informative, with $a_\beta = b_\beta = 2$, $m_0 = \frac{1}{2}\bm{1}$,
$\Sigma_0 = \frac{1}{K}\bm{I}_d$, $\nu_0 = 10^{-2}$, $\kappa_0 = 5$; given the
huge volume of data, we do not spend much effort tuning these.
We also use a hot-start initialization of $(\bm{z}, \bm{\mu}, \bm{\Sigma})$
using parameters obtained from a run of $K$-Means, while $\bm{\beta}$ is
initialized to $\bm{1}$.



\end{document}
